\documentclass[a4paper, 12pt]{article}

\usepackage[utf8]{inputenc}
\usepackage[T1]{fontenc}
\usepackage{textcomp}
\usepackage{amssymb}
\usepackage{newtxtext} \usepackage{newtxmath}
\usepackage{amsmath, amssymb}
\newtheorem{problem}{Problem}
\newtheorem{example}{Example}
\newtheorem{lemma}{Lemma}
\newtheorem{theorem}{Theorem}
\newtheorem{problem}{Problem}
\newtheorem{example}{Example} \newtheorem{definition}{Definition}
\newtheorem{lemma}{Lemma}
\newtheorem{theorem}{Theorem}
\DeclareMathAlphabet{\mathcal}{OMS}{cmsy}{m}{n}

\begin{document}

$(1)$ El concepto: ¿qué sucede si situamos un hombre verdaderamente bueno en un
mundo reinado por la vanidad, el egoísmo, la apariencia, y la frivolidad?

\begin{itemize}
    \item La novela como un \textit{ensayo}, o experimento: utilizar narrativa y
        personajes para explorar ideas abstractas. 
    \item Tolstoi y su optimismo respecto a la acción del hombre bueno.
\end{itemize}

$(2)$ La trama, con énfasis en el encuentro entre Rogozhin y Myshkin. 

\begin{itemize}
    \item Myshkin como Cristo.
    \item Rogozhin y Myshkin como dos caras de la misma moneda.
    \item La dualidad de Cristo como algo borrado del cristianismo: Jung,
        \textit{Aion}; Pistis Sophia.
    \item Nastasya como mujer caída. La contradicción entre rebelarse contra una
        sociedad que la condena y, al mismo tiempo, internalizar dicha condena.
\end{itemize}

~

$(3)$ Myshkin como el \textit{Poor knight} de Pushkin y como una interpretación
seria del Quijote.

\begin{itemize}
    \item El Quijote, visto desde fuera, es cómico, pero desde su perspectiva el
        mundo es serio y eminentemente trágico.
    \item El Quijote, por mucho que lo adoramos, indudablemente empeora el mundo
        en el que vive: libera prisioneros, perjudica a quienes quiere salvar,
        pone constantemente en peligro a Sancho Panza.
    \item Si bien Cervantes dice que don Quijote muere una muerte «sosegada y
        cristiana», el fin del Quijote es trágico: todo ha sido un sueño, todo
        ha sido en vano, deja atrás un mundo igual o peor a aquel que le vio
        nacer.
\end{itemize}








\end{document}



